% slots for James to write into
% red todo text shows in the pdf and % lines are notes that never print

\subsection*{Methodology}

\todo{Which traces: primary is the 1,034 traces of Qwen3-4B thinking, seed 42 (the Task 1 reasoning run and the Task 3 base). Replication on 10 more thinking runs (4B seeds 66, 73, 137, 255; 1.7B and 8B seeds 42, 66, 73), 11,374 traces in total. Correctness is EX, since Task 1 showed EM misjudges these models. Every trace is included, unfinished ones too.}

\paragraph{Features.}
\todo{For each feature give its definition, why it is relevant to Text-to-SQL, and how it was computed (Table~\ref{tab:features}). Relevance: length tracks planning effort on joins and nesting but can also mean the model got lost; ``alternatively'' marks weighing another formulation; ``wait'' marks self checking; drafts measure drafting and redrafting SQL; answer in trace is a surface faithfulness check \citep{chen2025reasoning}.}

\begin{table}[htbp]
\centering
\caption{Trace features: definition and computation, and why each matters for Text-to-SQL.}
\label{tab:features}
\scriptsize
\begin{tabularx}{\linewidth}{@{}>{\raggedright\arraybackslash}p{0.17\linewidth}>{\raggedright\arraybackslash}X>{\raggedright\arraybackslash}p{0.33\linewidth}@{}}
\toprule
Feature & Definition and computation & Relevance to Text-to-SQL \\
\midrule
Log reasoning tokens & Natural log of the number of tokens before \texttt{</think>}, as counted by vLLM. Log because counts are right skewed (median 498, maximum 7,174) & Harder queries need more planning (joins, nesting, schema linking), but a long trace can also mean the model got lost \citep{tian2026week3} \\
``Alternatively'' per 1,000 words & Whole-word count of ``alternatively'' in the lowercased trace, divided by the trace's word count, times 1,000. Words are runs of letters, digits and underscores, with the \texttt{<think>} tag removed & Marks exploring another formulation, such as a JOIN against a subquery or which table holds a column (``explore another branch'', \citealp{tian2026week3}) \\
``Wait'' per 1,000 words & The same computation for the word ``wait'' & Marks self checking a column name, condition or join key (``reconsider / verify'', \citealp{tian2026week3}) \\
SQL drafts & Count of the uppercase whole word SELECT in the trace. Case sensitive, since drafted SQL is in capitals and lowercase ``select'' is usually prose & A SQL specific measure of drafting and redrafting, including finding the answer and then ``correcting'' it away \citep{tian2026week3} \\
Answer in trace (0/1) & 1 if the final SQL appears verbatim in the trace after normalizing both (lowercase, collapsed whitespace, no spaces around punctuation, one quote style, no trailing semicolon) & A surface check on faithfulness: whether the visible reasoning contains the answer given \citep{chen2025reasoning} \\
\bottomrule
\end{tabularx}

\vspace{2pt}
\begin{minipage}{\linewidth}\scriptsize
\textit{Note.} A query with a subquery contains two SELECTs, so drafts slightly over count. Word counts are crude lexical measures. The first four features get the full set of tests; answer in trace gets a single test
\end{minipage}
\end{table}

\paragraph{Statistical tests.}
\todo{Name each test and why it fits the data: Mann--Whitney U with rank-biserial effect size for feature vs correctness (skewed, often zero, no normality assumption); Kruskal--Wallis and Spearman for feature vs difficulty (rank based, Spearman uses the ordering); AUROC overall and within each difficulty (threshold free, within level removes the confound that hard questions are both long and often wrong); logistic regression of correctness on difficulty plus the feature, then plus log length, with errors clustered by gold query (does a feature add anything beyond length; paraphrase pairs are not independent, \citealp{miller2024errorbars}); Fisher's exact test for answer in trace vs correctness (small cell); replication on 10 more runs as the safeguard for unadjusted exploratory tests.}

\begin{table}[htbp]
\centering
\caption{Statistical tests and why each fits the data.}
\label{tab:tests}
\scriptsize
\begin{tabularx}{\linewidth}{@{}>{\raggedright\arraybackslash}p{0.2\linewidth}>{\raggedright\arraybackslash}p{0.3\linewidth}>{\raggedright\arraybackslash}X@{}}
\toprule
Relates & Test & Why it fits \\
\midrule
Feature and correctness & Mann--Whitney U with rank-biserial effect size & Features are skewed and often zero, so a rank test avoids assuming normality \\
Feature and difficulty & Kruskal--Wallis over the four levels; Spearman $\rho$ with difficulty coded 0 to 3 & Rank based; Spearman also uses the ordering of difficulty \\
Separating right from wrong & AUROC overall and within each difficulty level & Threshold free; within each level it removes the confound that hard questions are both longer and more often wrong \\
Correctness with difficulty held fixed & Logistic regression of correctness on difficulty and the feature, then also log length; standard errors clustered by gold query & Tests the association at fixed difficulty, then whether a feature adds anything beyond length; clustering because paraphrase pairs share a gold query and are not independent \citep{miller2024errorbars} \\
Answer in trace and correctness & Fisher's exact test on the 2 $\times$ 2 table & One cell is small (about 50 traces without the answer) \\
Robustness & Key results recomputed on the 10 replication runs & Results from a single run can be unstable \citep{tian2026week3} \\
\bottomrule
\end{tabularx}

\vspace{2pt}
\begin{minipage}{\linewidth}\scriptsize
\textit{Note.} The tests are exploratory and $p$ values are not adjusted for multiple testing; replication across runs is the safeguard
\end{minipage}
\end{table}

\subsection*{Results}

\paragraph{Descriptive statistics.}
\todo{Present Table~\ref{tab:descriptives}. Answer in trace: 95.2\% of final queries appear in the trace (98.4\% easy to 89.8\% extra); accuracy 79.1\% when in the trace and 86.0\% when not.}

\begin{table}[htbp]
\centering
\caption{Trace features for Qwen3-4B thinking (seed 42) by correctness and difficulty, median [interquartile range].}
\label{tab:descriptives}
\scriptsize
\setlength{\tabcolsep}{3pt}
\resizebox{\linewidth}{!}{\begin{tabular}{@{}ll ccccc@{}}
\toprule
Feature & Group & Easy & Medium & Hard & Extra & All \\
\midrule
$n$ & right / wrong & 228 / 20 & 375 / 71 & 121 / 53 & 97 / 69 & 821 / 213 \\
\midrule
Reasoning tokens & right & 290.5 [231, 398] & 385 [289.5, 570.5] & 798 [543, 1,288] & 928 [586, 1,492] & 435 [289, 729] \\
 & wrong & 550.5 [309.8, 1,342.5] & 896 [529, 1,411.5] & 1,259 [536, 2,238] & 1,419 [929, 2,509] & 1,082 [575, 1,929] \\
\addlinespace
``Alternatively'' per 1,000 words & right & 0 [0, 0] & 0 [0, 1.9] & 2.3 [0, 4.0] & 2.6 [0, 3.8] & 0 [0, 2.7] \\
 & wrong & 2.7 [0, 3.9] & 1.3 [0, 3.4] & 2.6 [1.2, 3.8] & 2.2 [1.2, 3.7] & 2.3 [0, 3.8] \\
\addlinespace
SQL drafts (SELECT count) & right & 1 [1, 2] & 2 [1, 2] & 4 [2, 8] & 4 [2, 8] & 2 [1, 3] \\
 & wrong & 2 [2, 4] & 2 [2, 4.5] & 6 [2, 14] & 5 [2, 12] & 4 [2, 9] \\
\addlinespace
``Wait'' per 1,000 words & right & 6.2 [4.9, 8.9] & 5.7 [4.2, 8.1] & 5.6 [4.0, 7.2] & 5.4 [3.5, 6.8] & 5.8 [4.2, 8] \\
 & wrong & 4.7 [3.8, 8.5] & 6.2 [4.1, 7.9] & 5.1 [3.8, 7.1] & 5.1 [3.2, 7.8] & 5.2 [3.8, 7.7] \\
\bottomrule
\end{tabular}
}

\vspace{2pt}
\begin{minipage}{\linewidth}\scriptsize
\textit{Note.} Raw token counts are shown for readability; the tests use their log. Easy questions have only 20 wrong answers.
\end{minipage}
\end{table}

\paragraph{Statistical results.}
\todo{Present Table~\ref{tab:singletests}, Table~\ref{tab:regression} and Table~\ref{tab:replication}. Fisher's exact test for answer in trace vs correct: $p$ = 0.285.}
% length mann whitney p 1.9e-35 rank biserial 0.55, auroc 0.78 overall and 0.61 to 0.79 within levels
% with difficulty fixed each 2.7 fold longer trace multiplies the odds of being right by 0.33
% alternatively and drafts matter alone (or 0.852 p 8.9e-5 and 0.953 p 0.002) but not beyond length (p 0.48 and 0.22)
% wait never matters (auroc 0.45, p 0.59, the mann whitney p 0.037 is a negligible effect)
% replication length p under 1e-8 in every run, alternatively and wait never significant beyond length, drafts only once (1.7b seed 73)

\begin{table}[htbp]
\centering
\caption{Single feature tests against correctness and difficulty.}
\label{tab:singletests}
\scriptsize
\setlength{\tabcolsep}{3pt}
\begin{tabular}{@{}l rr rr rr rrrrr@{}}
\toprule
 & \multicolumn{2}{c}{Mann--Whitney} & \multicolumn{2}{c}{Kruskal--Wallis} & \multicolumn{2}{c}{Spearman $\rho$} & \multicolumn{5}{c}{AUROC, wrong vs right} \\
\cmidrule(lr){2-3}\cmidrule(lr){4-5}\cmidrule(lr){6-7}\cmidrule(lr){8-12}
Feature & $p$ & $r_{rb}$ & $H$ & $p$ & Difficulty & Length & Easy & Med. & Hard & Extra & All \\
\midrule
Log reasoning tokens & $1.9\times10^{-35}$ & 0.55 & 353 & $3.6\times10^{-76}$ & 0.58 & 1.00 & 0.75 & 0.79 & 0.61 & 0.65 & 0.78 \\
``Alternatively'' per 1,000 words & $6.0\times10^{-14}$ & 0.30 & 134 & $6.7\times10^{-29}$ & 0.33 & 0.60 & 0.69 & 0.65 & 0.52 & 0.51 & 0.65 \\
SQL drafts (SELECT count) & $8.3\times10^{-19}$ & 0.38 & 235 & $1.2\times10^{-50}$ & 0.45 & 0.73 & 0.75 & 0.69 & 0.57 & 0.56 & 0.69 \\
``Wait'' per 1,000 words & 0.037 & -0.09 & 23 & $4.4\times10^{-5}$ & -0.15 & -0.12 & 0.38 & 0.51 & 0.45 & 0.50 & 0.45 \\
\bottomrule
\end{tabular}


\vspace{2pt}
\begin{minipage}{\linewidth}\scriptsize
\textit{Note.} AUROC above 0.5 means higher values go with wrong answers. Rank-biserial $r_{rb}$ = 2 $\times$ AUROC $-$ 1. Spearman $\rho$ is with difficulty as an ordered scale (0 to 3) and with log length.
\end{minipage}
\end{table}

\begin{table}[htbp]
\centering
\caption{Logistic regression of correctness on each feature, odds ratio [95\% CI].}
\label{tab:regression}
\scriptsize
\begin{tabular}{@{}l ll@{}}
\toprule
Feature & Difficulty + feature & Difficulty + log length + feature \\
\midrule
Log reasoning tokens & 0.325 [0.247, 0.428], $p$ = $1.4\times10^{-15}$ &  \\
``Alternatively'' per 1,000 words & 0.852 [0.786, 0.923], $p$ = $8.9\times10^{-5}$ & 0.968 [0.885, 1.059], $p$ = 0.482 \\
SQL drafts (SELECT count) & 0.953 [0.924, 0.983], $p$ = 0.002 & 1.016 [0.991, 1.042], $p$ = 0.221 \\
``Wait'' per 1,000 words & 1.015 [0.962, 1.070], $p$ = 0.589 & 0.977 [0.926, 1.032], $p$ = 0.404 \\
\bottomrule
\end{tabular}


\vspace{2pt}
\begin{minipage}{\linewidth}\scriptsize
\textit{Note.} Every model includes difficulty. Odds ratios are per one unit of the feature; for log length one unit is about 2.7 times longer. Standard errors are clustered by gold query.
\end{minipage}
\end{table}

\begin{table}[htbp]
\centering
\caption{The main Task 2 results recomputed on every thinking run.}
\label{tab:replication}
\scriptsize
\setlength{\tabcolsep}{3pt}
\begin{tabular}{@{}l rr rr rrr rrr@{}}
\toprule
 & & & \multicolumn{2}{c}{Median tokens} & \multicolumn{3}{c}{Log length} & \multicolumn{3}{c}{$p$ beyond length} \\
\cmidrule(lr){4-5}\cmidrule(lr){6-8}\cmidrule(lr){9-11}
Run & EX & Unfinished & Right & Wrong & AUROC & OR & $p$ & Altern. & Drafts & Wait \\
\midrule
4B, seed 42 & 79.4 & 0 & 435 & 1,082 & 0.776 & 0.325 & $1.4\times10^{-15}$ & 0.48 & 0.221 & 0.40 \\
4B, seed 66 & 78.7 & 2 & 425 & 991.5 & 0.765 & 0.366 & $5.5\times10^{-14}$ & 0.54 & 0.085 & 0.82 \\
4B, seed 73 & 79.0 & 0 & 431 & 1,067 & 0.780 & 0.287 & $1.4\times10^{-18}$ & 0.75 & 0.112 & 0.27 \\
4B, seed 137 & 78.9 & 1 & 441.5 & 1,013 & 0.768 & 0.326 & $2.6\times10^{-16}$ & 0.53 & 0.067 & 0.25 \\
4B, seed 255 & 78.7 & 0 & 436.5 & 1,087.5 & 0.767 & 0.306 & $6.6\times10^{-17}$ & 0.46 & 0.220 & 0.91 \\
1.7B, seed 42 & 71.8 & 2 & 439.5 & 903 & 0.743 & 0.425 & $2.3\times10^{-10}$ & 0.46 & 0.076 & 0.08 \\
1.7B, seed 66 & 72.0 & 1 & 426 & 937.5 & 0.737 & 0.434 & $2.5\times10^{-9}$ & 0.49 & 0.218 & 0.23 \\
1.7B, seed 73 & 72.3 & 3 & 441.5 & 976.5 & 0.747 & 0.430 & $7.1\times10^{-10}$ & 0.20 & $4.8\times10^{-4}$ & 0.26 \\
8B, seed 42 & 78.9 & 2 & 509 & 1,252 & 0.770 & 0.331 & $2.1\times10^{-16}$ & 0.58 & 0.532 & 0.37 \\
8B, seed 66 & 79.1 & 3 & 514 & 1,266 & 0.777 & 0.302 & $9.9\times10^{-17}$ & 0.67 & 0.324 & 0.35 \\
8B, seed 73 & 79.9 & 2 & 513 & 1,215 & 0.766 & 0.337 & $8.0\times10^{-15}$ & 0.37 & 0.183 & 0.07 \\
\bottomrule
\end{tabular}


\vspace{2pt}
\begin{minipage}{\linewidth}\scriptsize
\textit{Note.} ``Unfinished'' counts traces that reached the 8,192-token cap. Length OR and $p$ come from the logistic model with difficulty. The last three columns are each feature's $p$ value when added to difficulty and log length.
\end{minipage}
\end{table}

\paragraph{Visualizations.}
\todo{Present Figure~\ref{fig:boxplots} (required) and Figure~\ref{fig:lengthquartile}.}
% hard is the exception in figure 6, shortest quarter 68 percent vs 86 for the second, then 72 and 52

\begin{figure}[htbp]
\centering
\includegraphics[width=\linewidth]{boxplots_full-qwen3-4b-think.png}
\caption{Trace features by difficulty and correctness for Qwen3-4B thinking, seed 42 (blue right, orange wrong, outliers hidden).}
\label{fig:boxplots}
\end{figure}

\begin{figure}[htbp]
\centering
\includegraphics[width=0.65\linewidth]{accuracy_by_length_full-qwen3-4b-think.png}
\caption{EX by trace length quartile within each difficulty level. Each point is a quarter of that level's questions, placed at its median length.}
\label{fig:lengthquartile}
\end{figure}

\paragraph{Answering RQ2.}
\todo{Interpret with numbers: length dominates; length also tracks difficulty; hard questions are the exception; branching and drafting are proxies for length; ``wait'' carries no signal; faithfulness check shows no relation; difference from the lecture notes; bridge to Task 3 (shortest trace tie break); limitations.}
% one sentence answer from the outline
% length is the only feature that relates reliably to correctness, wrong traces are about 2.5 times longer at every difficulty and across 11 runs and three sizes
